\documentclass[12pt, a4paper]{article}

% Paquetes requeridos
\usepackage[utf8]{inputenc}
\usepackage[T1]{fontenc}
\usepackage[spanish]{babel}
\usepackage{geometry}
\usepackage{amsmath}
\usepackage[american]{circuitikz}
\usetikzlibrary{babel} % Evita conflictos entre las flechas de circuitikz y babel-spanish

% Configuracion de la pagina
\geometry{a4paper, margin=2.5cm}

\begin{document}

\section*{Solución del Ejemplo 4.1: Modelo con Fuente Dependiente}

\subsection*{Planteamiento del Problema}
Hallar la expresión analítica para la tensión de salida $V_{out}$ en el circuito de la Figura \ref{fig:sol-ejm-4-1}, basándose en el análisis de mallas y de nodos.

\begin{figure}[h!]
    \centering
    \begin{circuitikz}[american, scale=1.2, transform shape]
        
        % Raíl inferior (Nodo común 'b')
        \draw (0,0) -- (7.5,0) node[ocirc] {}; 
        \node[below] at (3.5,0) {\textbf{b} (Referencia común)};

        % --- PUERTO DE ENTRADA ---
        \draw (0,0) to[V, l=$V_g$] (0,3) -- (2,3) node[circ] (A) {};
        \node[above] at (A) {a};
        \draw (A) to[R, l=$R_g$, i=$i_1$] (2,0);
        
        % --- PUERTO DE SALIDA ---
        \node[circ] (C) at (4,3) {};
        \node[above] at (C) {c};
        \draw (C) to[cI, l=$\beta i_1$] (4,0);
        \draw (C) -- (6,3) node[circ] {} to[R, l=$R_c$] (6,1.5) to[V, l=$V_{cc}$] (6,0);
        
        % Terminales de salida
        \draw (6,3) -- (7.5,3) node[ocirc] {};
        \draw[<->] (7.5,2.8) -- (7.5,0.2) node[midway, right] {$V_{out}$};

    \end{circuitikz}
    \caption{Circuito del Ejemplo 4-1.}
    \label{fig:sol-ejm-4-1}
\end{figure}

\subsection*{Desarrollo Matemático}

\textbf{Paso 1: Análisis del puerto de entrada.} \\
En el lado izquierdo del circuito, la resistencia $R_g$ se encuentra conectada en paralelo directamente a la fuente de tensión independiente $V_g$. Aplicando la Ley de Ohm, la corriente $i_1$ que desciende por $R_g$ es:
\begin{equation}
    i_1 = \frac{V_g}{R_g} \label{eq:i1}
\end{equation}

\textbf{Paso 2: Análisis del puerto de salida.} \\
En el lado derecho del circuito, tomamos el nodo inferior \textbf{b} como tierra o referencia ($V_b = 0\,\text{V}$). Dado que los terminales de salida están en circuito abierto, la tensión de salida $V_{out}$ corresponde exactamente a la tensión en el nodo \textbf{c} ($V_{out} = V_c$).

Aplicamos la Ley de Corrientes de Kirchhoff (LKC) en el nodo \textbf{c}. Como no fluye corriente hacia los terminales abiertos de la derecha, la sumatoria de las corrientes que salen del nodo \textbf{c} hacia las ramas inferiores debe ser igual a cero:
\begin{equation*}
    \beta i_1 + \frac{V_{out} - V_{cc}}{R_c} = 0
\end{equation*}

\textbf{Paso 3: Obtención de la expresión de $V_{out}$.} \\
Despejamos algebraicamente $V_{out}$ de la ecuación nodal anterior:
\begin{align*}
    \frac{V_{out} - V_{cc}}{R_c} &= -\beta i_1 \\
    V_{out} &= V_{cc} - \beta i_1 R_c \tag{2} \label{eq:vout_parcial}
\end{align*}

Finalmente, sustituimos la expresión de $i_1$ hallada en la Ecuación \eqref{eq:i1} dentro de la Ecuación \eqref{eq:vout_parcial}:
\begin{equation}
    V_{out} = V_{cc} - V_g \left( \beta \frac{R_c}{R_g} \right)
\end{equation}

\end{document}